% Conclusion chapter draft -- approx. 1.5 pages
% Placeholders in \emph{[brackets]} mark spots to fill with concrete numbers from the results tables/figures.

\section{Summary of contributions}

This thesis investigated the use of cell-free DNA (cfDNA) fragmentation patterns
for non-invasive cancer detection, with an emphasis on enlarging the feature
repertoire described by Zhou \emph{et al.} and on validating the resulting
classifier across multiple independent cohorts. A generalised
modelling pipeline was built around fourteen per-region or whole-genome
fragmentation features (length, FSD, FSR, PFE, coverage, ends, OCF, IFS, WPS,
EDM, iEDM, eEDM, eoEDM and ocEDM). Four of these features --- iEDM, eEDM,
eoEDM and ocEDM --- are introduced for the first time as extensions of the
end-motif diversity (EDM) family. Each feature is modelled independently
with a linear support-vector machine following per-fold standardisation,
sample-wise LOWESS GC correction (for the per-region features) and PCA dimensionality
reduction; the per-feature out-of-fold probabilities are then combined by a
second-layer SVM into an ensemble score.

\section{Per-feature and ensemble performance}

On the Cristiano \emph{et al.} cohort, the canonical
\texttt{svm\_linear\_pca150\_gc\_mapq30} configuration delivered AUCs of
\emph{[Zhou-comparable range]} for the nine features overlapping with the
Zhou panel, broadly reproducing the published baseline and confirming that
the modelling choices generalise across implementations. The four novel
features achieved AUCs of \emph{[insert range]}, with eoEDM and ocEDM in
particular reaching values competitive with the strongest per-region
features. The 14-feature ensemble exceeded every individual feature on the
Cristiano cross-validation, indicating that fragmentation signals captured by
distinct features are partially complementary rather than wholly redundant.

A stage-stratified analysis showed the expected monotone trend of
discriminability with disease stage, with the largest AUC gap between
healthy controls and stage~IV cancers and a substantially smaller separation
at stage~I. Cancer-type-stratified results highlighted that the per-feature
ranking is not invariant: features that lead overall do not necessarily
dominate every histology, motivating ensemble combination over single-feature
classifiers.

\section{External generalisation}

Generalisability was assessed by applying the Cristiano-trained models to two
internal cohorts (PreveLynch and GenoScan) and, conversely, by training
PreveLynch-based classifiers and evaluating them on the held-out PreveLynch
test split and on GenoScan. On the PreveLynch test split the
Cristiano-trained ensemble retained \emph{[insert AUC]} of its cross-validation
performance, demonstrating transferability across sequencing platforms and
sample preparations. Performance on GenoScan was \emph{[higher/lower than
PreveLynch]}, although the small cohort size (\(n=\emph{[insert N]}\))
limits the statistical resolution of that comparison. Sensitivity at
clinically relevant specificities (95\% and 85\%) followed the same
pattern as the AUCs and is reported per feature and per cohort in the
sensitivity tables.

\section{Limitations}

Several limitations of the present study should be highlighted. First, the
training cohorts are single-centre and modest in size relative to the
diversity of human cancers; the PreveLynch cohort in particular contains
primarily colorectal cancers and healthy controls, restricting the
cancer-type generality of LY-trained classifiers. Second, the GenoScan
external set is small, so reported AUCs there are noisier than for
PreveLynch. Third, fragmentation features were computed on a fixed open-chromatin
region set (\(\sim\)560k regions liftover'ed to GRCh37); changing the
reference region pool, the MAPQ filter, or the LOWESS bandwidth would
likely shift absolute scores, although ranking among features should be
robust. Fourth, all models were trained with default SVM hyperparameters
beyond kernel choice and PCA components; a more thorough hyperparameter
search was not performed because exploratory runs (see hyperparameter
section) indicated diminishing returns relative to the cost.

\section{Future directions}

The results suggest several productive avenues for subsequent work. The most
direct extension is validation on additional, ideally multi-centre, cohorts
covering a broader spectrum of cancer types and clinical stages, which would
both stress-test the present ensemble and enable cancer-type-specific
threshold calibration. A second direction is the integration of orthogonal
cfDNA modalities --- methylation patterns, somatic mutation signals, or
copy-number information --- into the existing fragmentation-feature
ensemble, building on evidence that fragmentation alone leaves a non-trivial
fraction of early-stage cancers undetected. Finally, deploying the
classifier in a clinical context will require calibration of the operating
threshold against the desired specificity (e.g.\ 95\% or 99\%), incorporation
of pre-analytical batch effects into the training procedure, and
prospective evaluation in a screening setting rather than the
case-control design used here.
